\documentclass[11pt]{article}
\usepackage[a4paper,margin=27mm]{geometry}
\usepackage{amsmath,amssymb,amsthm,hyperref}
\newtheorem{theorem}{Theorem}
\newtheorem{lemma}[theorem]{Lemma}
\newtheorem{corollary}[theorem]{Corollary}
\newcommand{\wt}{\operatorname{wt}}
\title{Prime-ary Walsh splitting and proper cells in even dimension}
\author{Companion to \texttt{walsh\_positive\_splitting.tex}}
\date{30 September 2026}
\begin{document}\maketitle
\noindent\textbf{Status.} Independently mathematically reviewed;
not formally verified. These are continuous-density
constructions. Neither theorem claims an improved finite-dice size bound.
The proper-cell variant was independently proposed by two collaborating
agents. The prime-ary variant increases the ratio of the total uniform
drift to the drift spent in fully active cells.

Fix a prime $q$, a primitive $q$th root $\zeta$, and a total-degree
cutoff $n$. Use formal generators $H$ and $Y_{i,a}$, where $i\in[n]$
and $a\in\mathbb F_q^*$, over $\mathbb Q(\zeta)$. Quotient the
associative algebra by all words containing two $Y$ letters with the
same first index $i$, and by all terms of degree above $n$. For
$b\in\mathbb F_q^n$, let $\rho_b(H)=H$ and
$\rho_b(Y_{i,a})=\zeta^{ab_i}Y_{i,a}$.
Starting with $W_0=\exp(H+\sum_{i,a}Y_{i,a})$, set
\[
 W_\ell=\prod_{t=0}^{q-1}\rho_{t b_\ell}(W_{\ell-1}),
\]
in the displayed order.
A nonzero monomial determines a frequency vector $v\in\mathbb F_q^n$:
its coordinate is $a$ if $Y_{i,a}$ occurs, and zero otherwise.

\begin{lemma}\label{filter}
If $B_\ell$ has rows $b_1,\ldots,b_\ell$, every log term with frequency
$v\ne0$ has degree at least $1+\wt(B_\ell v)$. The zero-frequency
terms equal $q^\ell H$.
\end{lemma}
\begin{proof}
At level zero, the assertion is immediate. In the BCH expansion for
$q$ factors, the linear sum multiplies frequency $v$ by
$\sum_{t=0}^{q-1}\zeta^{t b_\ell\cdot v}$, which is zero when the new
coordinate is nonzero. A nonlinear term is a bracket of $k\ge2$
components with disjoint first-index supports (otherwise it vanishes).
Its degree is at least $k$ plus the sum of the previous Hamming weights,
which is at least $k+\wt(B_{\ell-1}v)$. This allows the new required
weight to increase by one. Zero frequency can occur for a nonzero
monomial only when no $Y$ occurs, since each first index appears at most
once; all brackets composed only of $H$ vanish.
\end{proof}

\begin{theorem}\label{qpath}
For each fixed prime $q$, there are $q^{O_q(n)}$ equal rational cells
supporting a permutation-fair family of continuous densities
$f_i\in\{0,q\}$. The activity indicators on a uniformly selected cell
are independent with probability $1/q$. In particular, the total mass
contributed to each die by cells on which all $n$ dice are active is
exactly $q^{1-n}$.
\end{theorem}
\begin{proof}
Choose a matrix $B\in\mathbb F_q^{L\times n}$ of rank $n$ such that
all nonzero images have weight at least $n$, with $L=O_q(n)$.
For completeness, independent uniform rows give a binomial weight with
mean $L(1-1/q)$ for every nonzero input. Taking $L$ a sufficiently large
constant multiple of $n$, a Chernoff bound and the union bound over
$q^n-1$ inputs prove existence.
Lemma~\ref{filter} gives $\log W_L=q^L H$. Substitute
$H=\sum_iX_i$ and $Y_{i,a}=X_i$ in the reduced algebra. The substitution
respects the repeated-first-index ideal. A factor indexed by
$t\in\mathbb F_q^L$ has coefficient
\[
 1+\sum_{a\ne0}\zeta^{a(B^Tt)_i}
 =q\,\boldsymbol1_{(B^Tt)_i=0}
\]
on $X_i$. Thus all velocities are real, nonnegative, and rational.
After rescaling cell duration by $q^{-L}$, the signature is
$\exp(\sum_iX_i)$. Surjectivity of $B^T$ makes its values uniform on
$\mathbb F_q^n$, which proves the activity and mass assertions.
\end{proof}

\begin{theorem}\label{proper}
For each even $n\ge2$ and each fixed prime $q$, there is an exact fair
continuous family with $q^{O_q(n)}$ equal rational cells, with densities
in $\{0,q\}$, such that every cell has at least one inactive die.
\end{theorem}
\begin{proof}
Choose the rows of $B$ in the hyperplane $\boldsymbol1^\perp$, with
\[
 \ker B=\mathbb F_q\boldsymbol1,
 \qquad \wt(Bv)\ge n\quad(v\notin\ker B).
\]
The same random-row argument on the quotient gives $L=O_q(n)$.
The proof of Lemma~\ref{filter} applies without injectivity: it kills
every nonkernel frequency. A nonzero kernel frequency has every
coordinate nonzero, so any surviving term uses all $n$ distinct first
indices and has degree exactly $n$, with no $H$ letters. Hence
\[
 \log W_L=q^L H+E_n,
\]
where $E_n$ is homogeneous of degree $n$ and is central after truncation.
Apply a fixed phase $c\in\mathbb F_q^n$ with $\sum_i c_i\ne0$ before
substitution. The factors now have activity characterized by
$(B^Tt+c)_i=0$. Since every $B^Tt$ has coordinate sum zero, no factor
has all coordinates active. The log still has the form $q^LH+E'_n$.

Append the same list of exponential factors in reverse order.
Reversing a product of exponentials of degree-one elements multiplies
its degree-$d$ log term by $(-1)^{d+1}$: word reversal is an
antiautomorphism and acts by that sign on a homogeneous Lie polynomial.
Since $n$ is even, the reversed log is $q^LH-E'_n$. Centrality now
shows that the concatenation has log $2q^LH$ exactly. Normalize the
cell durations by $1/(2q^L)$. All cells remain proper.
\end{proof}

\begin{corollary}
For $q=2$ and even $n$, the final reversal in Theorem~\ref{proper}
can be omitted by choosing the first row of $B$ to be $\boldsymbol1$.
\end{corollary}
\begin{proof}
That row lies in $\boldsymbol1^\perp$ because $n$ is even. Add enough
independent random rows afterwards to obtain the same distance property.
Setting $H=0$ in the first symmetrization gives
$\exp(\sum_iY_i)\exp(-\sum_iY_i)=1$. Every later symmetrization, and
every fixed phase, preserves this identity. Thus every nonzero log error
contains at least one $H$. A surviving nonzero kernel frequency also
requires all $n$ distinct $Y_i$, so its degree exceeds the cutoff $n$.
\end{proof}

\noindent\textbf{Limits.} Theorem~\ref{proper} allows activity as large
as $n-1$, so naive recursive finiteization still gives a poor recurrence.
Theorem~\ref{qpath} supplies a drift ratio $q^{n-1}$, but a quantitative
positive absorption lemma is still required to remove its full cells.
The cell count measures geometric complexity of continuous densities;
it is not the face count of a finite equiprobable die.
\end{document}
